% Imporditud: tools/import_eksperimendid.py
% Allikas: Eesti füüsikaolümpiaadi eksperimendi ülesannete
%   kogu (Rael Kalda, 2023), ülesanne Ü6, lahendus L6.
\setAuthor{EFO žürii}
\setRound{piirkonnavoor}
\setYear{1999}
\setNumber{P E2}
\setDifficulty{1}
\setTopic{vedelike mehaanika}
\setVahendid{võrdse ruumalaga kehad, ümmargune pliiats, joonlaud. Ühe keha aine tihedus on 7,8 $\cdot$ $10^{3}$ kg/m$^3$}
\setPunktid{10}
\setAllikas{Eksperimendi kogu Ü6, lahendus lk 36}
\setLahendusseis{toores}

\prob{Keha tihedus}
Määrata keha aine tihedus.

\hint

\solu
Katse idee. Joonlauast ja pliiatsist kaalu valmistamine --- 1 p. Joonlaua esialgne tasakaalustamine --- 1 p. Kehade tasakaalustamine joonlauast kaaludel --- 1 p. Teooria ja valemi tuletamine --- 4 p. Kauguste $l_1$ ja $l_2$ mõõtmine --- 1 p. Aine tiheduse arvutamine --- 1 p. Arvutatud tulemuse mõõteviga $\Delta\rho$/$\rho \leq$ 10\% või on analüüsitud mõõtmise täpsust kvalitatiivselt --- 1 p.
\probend
